\documentclass[11pt,a4paper]{article}
\usepackage[utf8]{inputenc}
\usepackage[spanish]{babel}
\usepackage{graphicx}
\usepackage{geometry}
\usepackage{fancyhdr}
\usepackage{xcolor}
\usepackage{titlesec}
\usepackage{enumitem}

% Configuración de márgenes
\geometry{top=2.5cm, bottom=2.5cm, left=2.5cm, right=2.5cm}

% Definición de colores corporativos
\definecolor{sabanaBlue}{HTML}{1b3a4b}
\definecolor{darkGrey}{HTML}{2c3e50}
\definecolor{accentGrey}{HTML}{7f8c8d}
\definecolor{timeRed}{HTML}{c0392b}

% Configuración de encabezado y pie de página
\pagestyle{fancy}
\fancyhf{}
\lhead{\color{sabanaBlue}\bfseries S1 Métodos Estadísticos Avanzados}
\rhead{\color{accentGrey}Juan Nicolás Patiño Rodríguez}
\lfoot{\color{accentGrey}Guía del Exponente — Universidad de La Sabana}
\rfoot{\color{sabanaBlue}\bfseries Página \thepage}

% Títulos de sección personalizados
\titleformat{\section}
  {\color{sabanaBlue}\normalfont\large\bfseries}
  {\thesection}{1em}{}

\begin{document}

% Portada / Cabecera limpia
\begin{center}
  {\fontsize{16}{20}\selectfont \color{sabanaBlue}\bfseries GUÍA DEL EXPONENTE: PRESENTACIÓN EN VIVO (MÁX. 150 SEGUNDOS)}\\[0.2cm]
  {\fontsize{12}{15}\selectfont \color{darkGrey}\bfseries Unidad 2 — Impacto de Hábitos en el Rendimiento Académico}\\[0.1cm]
  {\color{accentGrey}\itshape Juan Nicolás Patiño Rodríguez — Maestría en Inteligencia Artificial}\\[0.5cm]
  \hrule
\end{center}

\vspace{0.3cm}

\section*{Instrucciones de Dirección y Ritmo}
\begin{itemize}[noitemsep,topsep=0pt]
  \item \textbf{Límite de Tiempo Máximo:} 2.5 minutos (150 segundos) estrictos.
  \item \textbf{Estilo de Locución:} Energético, riguroso, asertivo y directo al punto.
  \item \textbf{Postura en Video:} Mirada directa a la cámara. Apóyate en esta guía como apuntes visuales fuera de cámara.
\end{itemize}

\hrule

\vspace{0.5cm}

% --- SLIDE 1 ---
\noindent
\textbf{\large Diapositiva 1: Introducción y Pregunta de Investigación} \hfill \textcolor{timeRed}{\textbf{[0:00 - 0:28 | 28s]}}
\begin{description}[noitemsep,topsep=2pt]
  \item[Visual en Pantalla:] Título del informe, logotipos y resumen de la cohorte unificada de $N = 382$ estudiantes.
  \item[Guion de Voz (Speech):]
    «Saludos. Este análisis explora el impacto del estilo de vida en el rendimiento académico, usando una cohorte común de 382 estudiantes de matemáticas y lengua portuguesa. La pregunta que guía el estudio es: ¿cómo se relacionan el consumo de alcohol y la zona de residencia con las calificaciones y las inasistencias, y qué tan distinto es este patrón entre hombres y mujeres?»
\end{description}

\vspace{0.4cm}

% --- SLIDE 2 ---
\noindent
\textbf{\large Diapositiva 2: Rigor del Diseño Intra-Sujeto} \hfill \textcolor{timeRed}{\textbf{[0:28 - 0:50 | 22s]}}
\begin{description}[noitemsep,topsep=2pt]
  \item[Visual en Pantalla:] Diagrama del cruce determinista de 13 llaves y ventajas metodológicas del diseño intra-sujeto.
  \item[Guion de Voz (Speech):]
    «Al cruzar ambas asignaturas usando 13 llaves demográficas, establecemos un diseño intra-sujeto de alta robustez. Al evaluar a los mismos individuos en matemáticas y lenguaje, controlamos las variables sociodemográficas y familiares fijas, logrando aislar con total precisión el impacto de los hábitos individuales sobre el rendimiento cognitivo de la cohorte.»
\end{description}

\vspace{0.4cm}

% --- SLIDE 3 ---
\noindent
\textbf{\large Diapositiva 3: Consistencia y Resiliencia Académica} \hfill \textcolor{timeRed}{\textbf{[0:50 - 1:14 | 24s]}}
\begin{description}[noitemsep,topsep=2pt]
  \item[Visual en Pantalla:] \textit{correlacion\_cruzada.png} (Dispersión Jittered con línea intra-sujeto $r=0.48$ y notas sobre correlaciones y resiliencia).
  \item[Guion de Voz (Speech):]
    «En la correlación cruzada vemos que la habilidad lingüística y matemática se asocian positivamente con r igual a cero coma cuarenta y ocho. Sin embargo, el consumo de alcohol no actúa como un determinante absoluto: en matemáticas la correlación es prácticamente nula y quince estudiantes con alto consumo logran notas sobresalientes, evidenciando resiliencia individual.»
\end{description}

\vspace{0.4cm}

% --- SLIDE 4 ---
\noindent
\textbf{\large Diapositiva 4: Efectos de Interacción de Género y Estratos} \hfill \textcolor{timeRed}{\textbf{[1:14 - 1:36 | 22s]}}
\begin{description}[noitemsep,topsep=2pt]
  \item[Visual en Pantalla:] \textit{interaccion\_genero\_alcohol.png} (Curvas de medias en 3 estratos balanceados con barras $\pm 1$ SE).
  \item[Guion de Voz (Speech):]
    «Al agrupar el alcohol en tres estratos balanceados y calcular barras de error estándar, corregimos el artefacto de celdas pequeñas donde solo cuatro mujeres tenían consumo extremo. En portugués, los varones sufren un deterioro marcado al subir de estrato mientras las mujeres retienen su promedio; en matemáticas, el solapamiento de errores confirma que el alcohol no genera una brecha significativa en mujeres.»
\end{description}

\vspace{0.4cm}

% --- SLIDE 5 ---
\noindent
\textbf{\large Diapositiva 5: Análisis de Brecha Geográfica y Resiliencia} \hfill \textcolor{timeRed}{\textbf{[1:36 - 1:53 | 17s]}}
\begin{description}[noitemsep,topsep=2pt]
  \item[Visual en Pantalla:] \textit{brecha\_urbano\_rural.png} (Violín y boxplots híbridos con solapamiento distribucional).
  \item[Guion de Voz (Speech):]
    «En la brecha urbano-rural observamos un rezago medio de uno a uno punto cuatro puntos, asociado a un mayor tiempo de viaje en zonas rurales. No obstante, las curvas de violín demuestran un solapamiento masivo: múltiples estudiantes rurales alcanzan calificaciones sobresalientes, confirmando que la geografía impone retos logísticos pero no un techo cognitivo.»
\end{description}

\vspace{0.4cm}

% --- SLIDE 6 ---
\noindent
\textbf{\large Diapositiva 6: Detección Formal de Atípicos en Inasistencias} \hfill \textcolor{timeRed}{\textbf{[1:53 - 2:10 | 17s]}}
\begin{description}[noitemsep,topsep=2pt]
  \item[Visual en Pantalla:] \textit{atipicos\_inasistencias.png} (Boxplots por asignatura, línea punteada = límite superior de Tukey, notas de asimetría de impacto).
  \item[Guion de Voz (Speech):]
    «La regla de Tukey identifica atípicos severos en menos del cinco por ciento del alumnado. El hallazgo clave es su asimetría pedagógica: mientras en portugués las ausencias recurrentes castigan la nota en casi dos puntos, en matemáticas la media de los atípicos es idéntica a la del resto, sugiriendo dinámicas de estudio autónomo diferenciadas.»
\end{description}

\vspace{0.4cm}

% --- SLIDE 7 ---
\noindent
\textbf{\large Diapositiva 7: Conclusiones e Implicaciones de IA} \hfill \textcolor{timeRed}{\textbf{[2:10 - 2:30 | 20s]}}
\begin{description}[noitemsep,topsep=2pt]
  \item[Visual en Pantalla:] Conclusiones en tarjetas (Hábitos complejos y resiliencia, equidad territorial, intervención pedagógica asimétrica).
  \item[Guion de Voz (Speech):]
    «Para concluir: el consumo de alcohol muestra impactos heterogéneos y no deterministas, afectando principalmente el área verbal en varones. La brecha rural refleja fricción logística y no capacidad, y el ausentismo demanda atención focalizada por materia. Diseñemos modelos de IA justos y políticas educativas basadas en datos reales. Muchas gracias.»
\end{description}

\end{document}
